\section{Local continuation and product solutions}
\label{cross:section}

\subsection{The local metric extension theorem}

The common cap gives an initial isometry. To enlarge it, the global
argument will use source-generated harmonic coordinates to identify
neighborhoods of two interior points. In these coordinates we must
continue equality of the metrics across a point where equality is
already known on one side. The following theorem is the local input.

\begin{theorem}[Local metric extension]\label{local:extension}
Let $n\geq3$.  Let $(N_1,g_1)$ and $(N_2,g_2)$ be smooth compact
connected Riemannian manifolds with nonempty smooth boundary, with
Dirichlet Green functions $G_1,G_2$.  Identify two interior harmonic
coordinate neighborhoods with one neighborhood $\Omega\subset
\mathbb R^n$ of $0$.  Suppose an open $W\subset\Omega$ contains, near
$0$, one side of a smooth hypersurface through $0$, and suppose
\[
 g_1=g_2\text{ on }W,\qquad
 G_1(x,y)=G_2(x,y)\quad(x,y\in W,\ x\ne y).
\]
Suppose there are finitely many smooth pairs $u_j^a$ harmonic for
$g_j$ on $\Omega$, agreeing on $W$, whose first-side coordinate
Hessians at $0$ span
\[
 \{H\in\operatorname{Sym}^2(\mathbb R^n):g_1^{ij}(0)H_{ij}=0\}.
\]
Then the two coordinate metrics agree on a neighborhood of $0$.
\end{theorem}

We prove this theorem in Section~\ref{sec:bootstrap}. Its proof first
constructs a mixed Green kernel that is harmonic for the first metric in
$x$ and for the second metric in $y$: initially it is $G_1(x,y)$
when $y\in W$ and $G_2(x,y)$ when $x\in W$. For each prescribed
compact region at positive separation from $x=y$, a sufficiently small
spatial dilation allows the continuation constructed in this section to
reach that region. Section~\ref{sec:balls} then continues it to a family
of small spheres centered at $y$, under a temporary bound that controls
a normalized difference of the inverse metrics, together with finitely
many derivatives, by the current radius. Its flux transfers
first-metric harmonic functions to second-metric harmonic functions.
In particular, the resulting functional sends $u_1^a$ to $u_2^a$ and
reproduces the constant and coordinate functions.

The remaining task is uniform control as the spheres shrink.
Section~\ref{sec:angular} estimates the density representing the flux
functional. Under the same temporary metric bound, this estimate and exact
reproduction of affine functions give an $L^2$ bound for each paired
harmonic difference divided by the squared radius.
Section~\ref{sec:bootstrap} combines it with the spanning harmonic
Hessians to improve the temporary bound strictly. This improvement
allows the shrinking parameter to tend to zero with the spatial dilation
fixed, forcing the two metrics to be conformal near the contact point.
The common harmonic coordinates remove the remaining conformal factor. The geometric
continuation in this section provides existence away from the diagonal;
its estimates need not remain bounded as the separation tends to zero.

\subsection{Product solutions and extension from a cross}

We begin with a statement for general scalar elliptic operators. Let
$L_i$ be a smooth scalar second-order elliptic operator in a coordinate
domain, with real positive definite principal matrix $g_i^{-1}$.
Lower-order terms are allowed. A \emph{product solution} is a function
$F(x,y)$ satisfying
\[
 L_1^xF=0,\qquad L_2^yF=0.
\]
Such a function has ordinary unique continuation, since $L_1^x+L_2^y$
is elliptic on the product. We shall construct extensions explicitly;
unique continuation will identify extensions already constructed.

The starting data consist of two compatible product sets, with one
variable restricted to a smaller observed set in each piece. A common
basis for a solution space and its restriction to the smaller set will
construct the extension. Common-basis methods for separately harmonic
extension have a classical antecedent in \citet{Zeriahi1982}; singular
systems also enter quantitative Runge approximation
\cite{RulandSalo2019}. We give the variable-coefficient argument,
including its quantitative dependence.

\begin{lemma}[Extension from a cross]\label{cross:series}
Let $D_i$ be bounded connected coordinate domains and
$S_i\Subset D_i$ nonempty open sets.  Suppose a product solution $F$ is
given consistently on
\[
 (D_1\times S_2)\cup(S_1\times D_2),
\]
with the sum of its two $L^2$ norms at most $M$.
Let $O_i\Subset D_i$ be open sets such that every $L_i$-solution $w$
on $D_i$ satisfies
\begin{equation}\label{cross:interpolation-input}
 \|w\|_{L^2(O_i)}
 \le C\|w\|_{L^2(S_i)}^{\gamma_i}
          \|w\|_{L^2(D_i)}^{1-\gamma_i},
 \qquad 0<\gamma_i<1,\qquad \gamma_1+\gamma_2>1.
\end{equation}
Then $F$ extends as a product solution to $O_1\times O_2$, agreeing
with both given pieces on their intersections.  On every compact
subset of this product its $C^m$ norm is at most $C_mM$.
The constants are uniform when the geometric margins, ellipticity,
coefficient bounds, interpolation constants and positive exponent
margin are uniform.  Each specified output order uses only finitely
many coefficient bounds.
\end{lemma}

\begin{proof}
Let $\mathcal H_1(D_1)$ be the closed subspace of $L^2(D_1)$ consisting
of distributional $L_1$-solutions.  Restriction
$R:\mathcal H_1(D_1)\to L^2(S_1)$ is compact by interior elliptic
estimates and injective by unique continuation.  The spectral theorem
for $R^*R$ gives an orthonormal basis $(w_j)$ of
$\mathcal H_1(D_1)$ such that
\begin{equation}\label{cross:singular-values}
 (w_i,w_j)_{L^2(S_1)}=\mu_j^2\delta_{ij},
 \qquad 0<\mu_j\le1.
\end{equation}
For every $N$,
\begin{equation}\label{cross:singular-decay}
 \mu_j\le C_N(1+j)^{-N}.
\end{equation}
Indeed, multiply a solution by a cutoff equal to one near
$\overline{S_1}$ and compactly supported in $D_1$.  Its $H^k$ norm is
bounded by its $L^2(D_1)$ norm.  Fourier truncation in a containing cube
gives a rank $O(A^n)$ approximation with operator error $O(A^{-k})$.
The minimax characterization of singular values yields
$\mu_j\le C_kj^{-k/n}$.  All these constants have the stated uniformity.

For $y\in S_2$ expand $F(\cdot,y)$ on $D_1$:
\[
 F(x,y)=\sum_j w_j(x)b_j(y),\qquad
 \|b_j\|_{L^2(S_2)}\le M.
\]
Equation~\eqref{cross:singular-values} continues these coefficients to
$D_2$ by the formula
\begin{equation}\label{cross:coefficient-extension}
 b_j(y)=\mu_j^{-2}\int_{S_1}F(x,y)\overline{w_j(x)}\,\dd x,
 \qquad \|b_j\|_{L^2(D_2)}\le M\mu_j^{-1}.
\end{equation}
The integral solves $L_2b_j=0$ and agrees with the original coefficient
on $S_2$.  Applying \eqref{cross:interpolation-input} gives
\[
 \|w_j\|_{L^2(O_1)}\le C\mu_j^{\gamma_1},\qquad
 \|b_j\|_{L^2(O_2)}\le CM\mu_j^{\gamma_2-1}.
\]
Consequently the series converges absolutely in $L^2(O_1\times O_2)$:
its summands are bounded by $CM\mu_j^{\gamma_1+\gamma_2-1}$.
Choose $N$ in \eqref{cross:singular-decay} so that
$N(\gamma_1+\gamma_2-1)>1$.  The limit solves both equations
in distributions.  Interior estimates for their elliptic sum give
the asserted smoothness and bounds on smaller sets.

We verify agreement with the second given piece; no assertion that
an arbitrary local solution extends to $D_1$ is needed.  Let $z$ be
orthogonal to the closure of $R\mathcal H_1(D_1)$ in $L^2(S_1)$.
The function
\[
 y\longmapsto\int_{S_1}F(x,y)\overline{z(x)}\,\dd x
\]
solves the second equation and vanishes on $S_2$, so it vanishes on
connected $D_2$.  Thus each second-piece slice belongs to that closed
range.  Its expansion in the orthonormal basis $w_j/\mu_j$ is complete,
and its coefficients are $\mu_jb_j(y)$ by
\eqref{cross:coefficient-extension}.  It follows that the series is
$F$ on $S_1\times O_2$; the series converges there also by the positive
exponent $\gamma_2$.  On $O_1\times S_2$ it is the original expansion,
with positive exponent $\gamma_1$.  This proves both agreements.
Only upper bounds for the singular values were used, so uniformity
does not require continuous choices of singular bases or positive
lower bounds for their singular values.
\end{proof}

Figure~\ref{cross:figure} summarizes the sets and the two estimates in
Lemma~\ref{cross:series}.
\begin{figure}[htbp]
 \centering
 \begin{tikzpicture}[x=1cm,y=1cm,font=\small,>=Stealth,
  line cap=round,line join=round]
  \colorlet{crossink}{blue!45!black}
  \fill[crossink!12] (0,0.70) rectangle (4.2,1.30);
  \fill[crossink!12] (1.05,0) rectangle (1.65,3.35);
  \fill[crossink!25] (1.05,0.70) rectangle (1.65,1.30);
  \draw[black!45] (0,0) rectangle (4.2,3.35);
  \draw[crossink,thin] (0,0.70)--(4.2,0.70)
    (0,1.30)--(4.2,1.30) (1.05,0)--(1.05,3.35)
    (1.65,0)--(1.65,3.35);
  \draw[densely dashed,thick] (0.58,0.37) rectangle (3.66,2.85);
  \node[fill=white,inner sep=3pt] at (2.65,2.30) {$O_1\times O_2$};
  \node[text=crossink] at (2.95,1.00) {$D_1\times S_2$};
  \node[text=crossink,rotate=90] at (1.35,2.10) {$S_1\times D_2$};
  \node[anchor=north] at (2.10,-0.12) {$D_1$ (first variable)};
  \node[rotate=90,anchor=south] at (-0.15,1.68) {$D_2$ (second variable)};
  \draw[->,thick] (4.60,1.68)--(5.32,1.68);
  \node[anchor=west,align=left,text width=6.10cm] at (5.60,1.68) {
    Compatible solutions on the cross\\[6pt]
    $\displaystyle
       \|w\|_{O_i}\leq C\|w\|_{S_i}^{\gamma_i}
                              \|w\|_{D_i}^{1-\gamma_i}$\\[6pt]
    $\displaystyle 0<\gamma_i<1,\qquad \gamma_1+\gamma_2>1$\\[8pt]
    The cross series constructs a solution\\
    on the smaller product $O_1\times O_2$.
  };
\end{tikzpicture}

 \caption{Compatible product solutions are given on the shaded cross.
 The $L^2$ interpolation estimates for single-variable solutions,
 with $\gamma_1+\gamma_2>1$, make the
 singular-function series converge on $O_1\times O_2$.
 The drawing is schematic: each factor is a multidimensional domain,
 and containment of $S_i$ in $O_i$ is not required.}
 \label{cross:figure}
\end{figure}

\subsection{A geometric criterion for local extension}

We next build the interpolation estimates needed in
Lemma~\ref{cross:series}.  For a metric $g$, call a smooth real function
$\theta$ a \emph{strict weight} at a point where $\dd\theta\ne0$ if
\begin{equation}\label{cross:strict-weight}
 \operatorname{Hess}_g\theta(e,e)
  +\operatorname{Hess}_g\theta(v,v)>0
 \quad\text{for all }v\perp e,\ |v|_g=1,
 \qquad e=\frac{\nabla^g\theta}{|\nabla^g\theta|_g}.
\end{equation}
The strict scalar Carleman estimate gives, after shrinking the
coordinate neighborhood,
\begin{equation}\label{cross:carleman}
 \tau^3\|e^{\tau\theta}w\|_{L^2}^2
 +\tau\|\nabla(e^{\tau\theta}w)\|_{L^2}^2
 \le C\|e^{\tau\theta}Lw\|_{L^2}^2,
 \qquad \tau\ge\tau_0,
\end{equation}
for compactly supported $w$.
This is the elliptic strict-weight estimate of the classical
Carleman theory; see \citet[Chapter~XXVIII]{HormanderIV2009} and
\citet[Definition~8.3 and Theorem~9(a)]{KochTataru2005}.
Here the hypotheses can also be checked directly.  For $L=-\Delta_g$,
write $v=e^{\tau\theta}w$ and
\[
 e^{\tau\theta}Le^{-\tau\theta}=A_\tau+B_\tau,\qquad
 A_\tau=-\Delta_g-\tau^2|\dd\theta|_g^2,\quad
 B_\tau=\tau(2\nabla^g\theta\cdot\nabla+\Delta_g\theta).
\]
Here $A_\tau$ is self-adjoint and $B_\tau$ is skew-adjoint for
metric volume.  Integration by parts gives
\begin{align*}
 \|(A_\tau+B_\tau)v\|^2
 &=\|A_\tau v\|^2+\|B_\tau v\|^2
       +([A_\tau,B_\tau]v,v),\\
 ([A_\tau,B_\tau]v,v)
 &=4\tau\int\operatorname{Hess}_g\theta(\nabla v,\nabla\overline v)
   +4\tau^3\int\operatorname{Hess}_g\theta
                 (\nabla\theta,\nabla\theta)|v|^2
   -\tau\int(\Delta_g^2\theta)|v|^2.
\end{align*}
Choose a real constant $m$ strictly between the negative of the least
unit tangential Hessian value and the unit normal Hessian value.
After shrinking the patch, these inequalities have a positive margin.
Add $4m\tau(A_\tau v,v)$ to the commutator form.  Its gradient
coefficient becomes $4\tau(\operatorname{Hess}_g\theta+mg)$,
and its leading zero-order coefficient is
\[
 4\tau^3|\dd\theta|_g^2
       \big(\operatorname{Hess}_g\theta(e,e)-m\big)>c\tau^3.
\]
The tangential gradient form is positive.  Its mixed and possibly
negative normal terms cost at most
$C\tau\|\nabla_e v\|^2$, and
\[
 \|\nabla_e v\|^2
 \le C\tau^{-2}\|B_\tau v\|^2+C\|v\|^2.
\]
Finally
$|4m\tau(A_\tau v,v)|\le\tfrac12\|A_\tau v\|^2+
C\tau^2\|v\|^2$.
For large $\tau$, the squared parts absorb the normal derivative
cost and the positive zero-order term absorbs the remaining errors.
This yields \eqref{cross:carleman}.  Smooth first- and zero-order terms
in $L$ cost only $C(\|\nabla v\|^2+\tau^2\|v\|^2)$ and are absorbed
in the same way.  The proof gives uniform constants on fixed compact
sets of strict weights and bounded coefficients.

The existence criterion is a condition on a hypersurface in the
product.  Hessians in the following statement use the product
connection of $g_1$ and $g_2$.

\begin{proposition}[Product extension criterion]\label{cross:criterion}
Let $\rho$ be smooth near $z_0=(x_0,y_0)$, with $\rho(z_0)=0$ and
$\dd_x\rho,\dd_y\rho\ne0$.  Put
\[
 a_i=\frac{\nabla_i\rho}{|\nabla_i\rho|^2},\qquad H=\operatorname{Hess}\rho.
\]
Suppose that at $z_0$
\begin{equation}\label{cross:hessian-test}
 H((a_1,-a_2),(a_1,-a_2))
 +H((v_1,v_2),(v_1,v_2))>0
 \quad
 \left(v_i\perp a_i,\ |v_i|=|a_i|\right).
\end{equation}
Every product solution given on $\{\rho>0\}$ near $z_0$ extends to
a neighborhood of $z_0$, agreeing on a smaller part of that side.
In fact, the construction uses only two compact product sets in
$\{\rho>\epsilon\}$ for some $\epsilon>0$.
The neighborhoods and estimates on smaller neighborhoods are uniform
under small smooth perturbations preserving the strict hypotheses,
provided the data on those compact sets are bounded.
\end{proposition}

\begin{proof}
We first split the product Hessian condition into a strict weight in
each factor. Their sum will lie below a defining function for the given
side, with a quadratic gap. That gap places two intersecting product
sets in the given side; the weights then yield interpolation exponents
greater than one half, allowing Lemma~\ref{cross:series} to apply.

\paragraph{Splitting the weight between the factors.}
We seek symmetric forms $P_i$ satisfying
\[
 P_i(a_i,a_i)+P_i(v_i,v_i)>0,
 \qquad H+K\,\dd\rho^{\otimes2}>P_1\oplus P_2
\]
for some $K>0$, with the vectors $v_i$ as in
\eqref{cross:hessian-test}. After division by $|a_i|^2$, the first
inequality is the strict-weight condition for a function with gradient
$\nabla_i\rho$ and Hessian $P_i$. The second inequality will give the
quadratic gap. To obtain these forms by convex separation, the dual
tests are positive semidefinite forms $D$ whose diagonal blocks are
\[
 d_i a_i^{\otimes2}+S_i,\qquad
 d_i\ge0,\quad S_i\ge0\text{ on }a_i^\perp,
 \quad \tr S_i=d_i|a_i|^2.
\]
This follows by separating the open convex cone consisting of positive
definite forms plus the allowed $P_1\oplus P_2$; the individual dual
cones are generated by $a_i^{\otimes2}+v_i^{\otimes2}$.
Normalize $\tr D=1$.  To obtain a suitable $K$, it suffices by
compactness to prove $H:D>0$ on tests with $D\,\dd\rho=0$.

Represent such a $D$ as a covariance matrix.  The normal components
measured by the two partial covectors are opposite.  The diagonal-block
condition makes each normal component uncorrelated with its own
tangential component, hence with both tangential components.
Their variances coincide, giving $d_1=d_2=d>0$ and
\[
 D=d(a_1,-a_2)^{\otimes2}+D_\perp,
 \qquad \tr(D_\perp)_{ii}=d|a_i|^2.
\]
The case $d=0$ would force $D=0$.  Among nonnegative tangential forms
with these two fixed traces, every extreme point has rank one:
on an image of rank $r\ge2$, a nonzero symmetric perturbation preserves
the two traces because $r(r+1)/2>2$, and small perturbations of both
signs preserve positivity.  The rank-one tests are precisely those
in \eqref{cross:hessian-test}, multiplied by $d$.
The asserted positivity follows.  Compactness of the normalized tests
then supplies $K$, as required.

Prescribe gradients $\dd_i\rho$ and Hessians $P_i$ for smooth functions
$\theta_i$ vanishing at the respective points.  In small coordinates
centered there, Taylor expansion gives
\begin{equation}\label{cross:quadratic-gap}
 \theta_1(x)+\theta_2(y)
 \le \widetilde\rho(x,y)-c(|x|^2+|y|^2),
 \qquad \widetilde\rho=\rho+K\rho^2/2,
\end{equation}
with $c>0$.  Each $\theta_i$ satisfies
\eqref{cross:strict-weight}.  The positive side of $\widetilde\rho$
is the positive side of $\rho$ in this neighborhood.

\paragraph{Placing the cross and obtaining interpolation.}
Use coordinates $(s_i,z_i)$ with $s_i=\theta_i$. Fix a small
$\kappa>0$ so that $\phi(s_i)=s_i-\kappa s_i^2$ remains a strict weight.
Choose successively a small lateral radius $R$, a height
$l\ll cR^2$, and $e\ll\kappa l^2$.  Take
\[
 D_i=\{-l<s_i<l+4e,\ |z_i|<R\}.
\]
Use a cutoff equal to one on the inner lateral half and for
$-l+3e<s_i<l+2e$, supported in the cylinder and in
$-l+2e<s_i<l+3e$.
Let $S_i\Subset D_i$ include neighborhoods of its top and lateral
derivative supports, where respectively
\[
 s_i>l+e/2\quad\text{or}\quad |z_i|>R/3,
\]
and an inner top slice near $s_i=l+5e/4$.
The two closed product sets $\overline{D_1\times S_2}$ and
$\overline{S_1\times D_2}$ lie strictly in the given side.
For a top slice, the sum of the two $s$ coordinates is positive.
For a lateral slice, it is greater than $-2l$, which is dominated
by the quadratic gap in \eqref{cross:quadratic-gap}.
Making the sets slightly smaller if necessary gives a common
$\epsilon>0$ with $\rho>\epsilon$ on their closures.

For an $L_i$-solution normalized by $\|w\|_{L^2(D_i)}=1$, apply
\eqref{cross:carleman} to this cutoff times $w$.  Interior estimates
bound the bottom derivative terms using the large norm and the weight
$\phi(-l+3e)$.  The remaining derivative terms use
$\|w\|_{L^2(S_i)}$ and a weight at most $\phi(l+3e)$.
On a smaller inner cylinder beginning at $s_i=-e$, the weight is at
least $\phi(-e)$.  Balancing the two exponentials yields
\eqref{cross:interpolation-input} with
\begin{equation}\label{cross:exponent}
 \gamma_i=
 \frac{\phi(-e)-\phi(-l+3e)}{
       \phi(l+3e)-\phi(-l+3e)}>\frac12.
\end{equation}
Indeed at $e=0$ the quotient is $1/2+\kappa l/2$.
Values of the balancing parameter below $\tau_0$ are covered by
increasing the constant.  The output cylinder may be chosen to include
a connected tube from the origin to the indicated top slice.

\paragraph{Constructing and identifying the extension.}
Lemma~\ref{cross:series} now gives a product solution on a neighborhood
of $z_0$. It agrees with the old solution there on the positive side:
from a sufficiently near positive-side point, increase $s_2$ into
the top slice.  Since $\partial_{s_2}\rho>0$ locally, the segment stays
on that side and in the larger constructed cylinder.  Agreement on
the slice and unique continuation along a neighborhood of the segment
prove agreement at the original point.
All choices have strict margins.  Keeping their finitely many compact
sets and shrinking their output neighborhoods once proves the stated
perturbation uniformity by \eqref{cross:carleman},
Lemma~\ref{cross:series} and interior elliptic estimates.
\end{proof}

\subsection{Gluing along a compact family of surfaces}

Proposition~\ref{cross:criterion} provides a germ across one surface
point.  The following formulation records the geometry needed to
combine these germs.  In particular, agreement is proved on specified
overlap components, rather than inferred from the existence of a cover.

\begin{lemma}[Propagation through a compact cylinder]
\label{cross:propagation}
Suppose a region in a product patch has smooth coordinates $(z,a)$
near $B\times[a_-,a_+]$, where $B$ is a compact base, possibly with
boundary or corners.  A product solution is given on a neighborhood
of the top and lateral boundary and on an open set $\mathcal S$.
Assume that $\mathcal S$ is upward closed in $a$ at fixed $z$ and
contains these boundary neighborhoods.  If every level $a=\text{constant}$
satisfies Proposition~\ref{cross:criterion} outside $\mathcal S$, with
the positive side pointing toward increasing $a$, the solution extends
to a neighborhood of the cylinder, agreeing on $\mathcal S$.
For fixed coordinates, compact sets and strict reference inequalities,
the extension has uniform compact-subset bounds under small smooth
perturbations, in terms of bounds on finitely many compact subsets of
the initially given region.
\end{lemma}

\begin{proof}
Starting from the top, let $a_*$ be the infimum of levels down to which
an extension agreeing with the seed has been obtained.  In the uniform
version include uniform bounds in this property.  At a non-seed point
of $B\times\{a_*\}$, Proposition~\ref{cross:criterion} uses compact
data strictly above that level.  Take a finite cover of the level by
such output neighborhoods and seed neighborhoods.  Their data are
contained above $a_*+\eta$ for one $\eta>0$, so an already reached
level supplies them with the needed bounds.

Shrink the outputs in $(z,a)$ coordinates to boxes
$B_i\times(a_*-d_i,a_*+d_i)$.
Every connected component of an overlap contains vertical segments
to a common height greater than $a_*$.  Both germs equal the preceding
extension there.  Unique continuation for $L_1^x+L_2^y$ makes them
identical on that overlap component.  The same argument gives agreement
with $\mathcal S$, because a vertical segment moving upward from a
seed point stays in the seed.  Boundary neighborhoods use the given
solution.  A finite subcover has a positive minimum output width and
therefore passes the putative last level, or includes the endpoint
$a_-$.  This proves continuation on the whole cylinder.
The same finite constructions retain all strict margins under small
perturbations.  Their Carleman, series and interior estimates prove
the uniform assertion.  Data bounded up to a limiting surface are
never required: each local construction uses compact sets strictly
above it.
\end{proof}

\subsection{Initialization at a fixed distance from the diagonal}

We apply the criterion to a mixed Green kernel.  Let $g_1,g_2$ be smooth
metrics in coordinate patches centered at $0$, obtained from interior
patches of compact Dirichlet manifolds.  Suppose they agree on an open
set $W$ containing the lower side of a smooth hypersurface through $0$.
After a linear coordinate change, assume
\[
 g_1(0)=g_2(0)=I,\qquad
 \{x_n<q(x')\}\subset W\text{ near }0,\qquad
 q(0)=0,\quad \dd q(0)=0,
\]
for a smooth function $q$.  No regularity of the rest of $\partial W$
is assumed.
Let $G_j$ be the Dirichlet Green kernels, normalized with metric volume,
and suppose $G_1=G_2$ on $W\times W$.
For a small dilation parameter $\lambda>0$, put
\begin{equation}\label{cross:dilation}
 g_{j,\lambda}(z)=g_j(\lambda z),\qquad
 G_{j,\lambda}(x,y)=\lambda^{n-2}G_j(\lambda x,\lambda y),
 \qquad W_\lambda=\lambda^{-1}W.
\end{equation}
On each fixed bounded patch the metrics tend smoothly to the Euclidean
metric, and $W_\lambda$ contains every fixed compact subset of
$\{z_n<0\}$ for sufficiently small $\lambda$.
Initially define, off the diagonal,
\begin{equation}\label{cross:mixed-data}
 \mathcal G_\lambda(x,y)=
 \begin{cases}
 G_{1,\lambda}(x,y),&y\in W_\lambda,\\
 G_{2,\lambda}(x,y),&x\in W_\lambda.
 \end{cases}
\end{equation}
The pieces fit and solve the first equation in $x$ and the second in $y$.

\begin{proposition}[Fixed-separation initialization]\label{cross:initial}
In this setting, let $n\ge3$ and fix $B>0$, $d>0$ and $\eta>0$.
For all sufficiently small $\lambda$, the data
\eqref{cross:mixed-data} have a product-solution continuation to a
neighborhood of
\[
 \{(x,y):|x|,|y|\le B,\ |x-y|\ge d\}.
\]
Every fixed $C^m$ norm on a smaller neighborhood is bounded uniformly
in $\lambda$.  The continuation agrees with $G_{1,\lambda}$ where
$y_n<-\eta$ and with $G_{2,\lambda}$ where $x_n<-\eta$, on the
corresponding overlaps.  The smallness threshold for $\lambda$ and
the constants may depend on $B,d,\eta,m$.
\end{proposition}

\begin{proof}
We first cross the limiting plane, then deform variable-radius tubes
to reach the prescribed separations.  All the sets and positive
separation margins used in the two steps are fixed before choosing
$\lambda$.

\paragraph{Crossing the plane and obtaining a common seed.}
We first extend slightly across the plane in one variable while keeping
the other at a comparable positive separation. The first interpolation
exponent can be made close to one; this compensates for the fixed
positive exponent obtained by propagation in the separated second
variable.

Fix a point $x_0$ of height zero and a separation scale $s>0$.
In the second factor take
\[
 D_2=\{s/2<|y-x_0|<2s\}
\]
and an observed ball $S_2$ about $x_0-se_n$.
In the first take a ball of radius $3\tau_0s$ centered at
$x_0-\tau_0se_n$, with $\tau_0>0$ fixed small enough that the two
large domains are separated.  Its observed ball has radius
$(1-\alpha)\tau_0s$, with $\alpha>0$ to be chosen.
These observed balls lie strictly in the lower half-space.

There is a two-constants estimate from $S_2$ to a slightly enlarged
closed annulus $4s/5\le|y-x_0|\le6s/5$, with an exponent
$\gamma_2>0$ independent of $\alpha$.  Here is a direct way to obtain
it with room for perturbations.  On a Euclidean annulus the weight
$|z-z_*|^{-1}$ is strict: its radial Hessian is
$2|z-z_*|^{-3}$ and each tangential Hessian value is
$-|z-z_*|^{-3}$.  Radial cutoffs in
\eqref{cross:carleman}, followed by exponential balancing, propagate
smallness from an inner ball to an intermediate ball using the norm
on an outer ball.  A finite chain of overlapping balls with enlarged
balls inside $D_2$ connects $S_2$ to the target annulus.  The annulus
is connected since $n\ge3$.  Multiplying the finitely many positive
exponents gives the asserted $\gamma_2$.  The same weights and margins
work for sufficiently small perturbations of the Euclidean metric.

In the first factor use this radial argument centered at
$x_0-\tau_0se_n$.  Keep the outer cutoff a fixed distance beyond radius
$\tau_0s$, and put the inner cutoff just inside radius
$(1-\alpha)\tau_0s$.  Take a concentric target ball of radius
$(\tau_0+c)s$, with $c>0$ small; it contains the observed ball and
a ball of radius $cs$ about $x_0$.
As $\alpha$ and $c$ decrease, the target weight approaches the inner weight.
The interpolation exponent therefore tends to one.  Choose these
constants so that $\gamma_1+\gamma_2>1$.
Lemma~\ref{cross:series} extends the kernel to a small ball about
$x_0$ times the target annulus.  On these fixed separated sets the
rescaled Green kernels and their derivatives are uniformly bounded:
the local Green singularity is $O(|x-y|^{2-n})$, and interior elliptic
estimates give derivative bounds away from the pole.

The extensions agree with both lower pieces.  Agreement for $y\in S_2$
comes from the first expansion.  For fixed $x$ below a sufficiently
small negative height, continue in $y$ through the connected target
annulus to compare with $G_{2,\lambda}$.  For a second point below that
height, compare with $G_{1,\lambda}$ by continuation in the first
variable from its observed ball.  All comparisons remain separated.

The local extensions must next be compared on overlaps. Retain much
smaller first balls and narrower annuli for coverage, keeping the larger
products as comparison domains. The unused margins allow both variables
to move into the known lower region. Overlapping retained products have
scales comparable to $|x-y|$. For two products of the same orientation,
lower their near-plane variable into the known region; the path stays
inside both larger first balls.  For products of opposite orientations,
both variables lie in retained near-plane balls. On the fixed compact
range of locations and positive separation scales, choose the retained
products with strict ball and annular margins, then take $\eta_0>0$
much smaller than the smallest retained width. Lower both variables by
\[
 h=\max(x_n,y_n,0)+2\eta_0.
\]
By making the retained balls a sufficiently small fraction of the larger
balls, the path stays in both larger products: its displacement is
smaller than their ball and annular margins. Its endpoint has both
heights below $-\eta_0$,
where both germs are the original kernel.  Unique continuation along
a neighborhood of the path proves agreement.  Choose
$\eta_0\le\eta$ and then $\lambda$ small enough for this truncated
lower half-space to lie in $W_\lambda$ on all sets used.

Interchanging the two variables and covering gives, for some fixed
$c_0>0$, a single-valued continuation in
\begin{equation}\label{cross:coarse-region}
 \min(x_n,y_n)<c_0|x-y|,
\end{equation}
on every fixed compact range of bounded locations and positive
separations.  The estimates are uniform in $\lambda$ on that range.
For coverage near the plane, take $x_0=(x',0)$ and
$s=|y-x_0|$ when the first variable has the smaller nonnegative height;
then $x$ is in the retained first ball and $y$ in the retained annulus
if $c_0$ is small.  Points already below the plane are covered by the
same construction near it or by the given cross farther below it.
The preceding overlap argument makes these choices compatible.

\paragraph{Sweeping tubes inward from the common seed.}
The region \eqref{cross:coarse-region} is now the seed for the second
stage. At fixed center $y$ and direction $\omega$, we move $x$ along the
ray from $y$: large radii lie in the seed, and decreasing the radius
will reach the target pairs. Choose $M>B+3$ and introduce
\begin{equation}\label{cross:tubes}
 t=M+y_n+e_0|y'|^2,\qquad r_a(y)=a t^{1+\sigma},\qquad
 x=y+r_a(y)\omega,\quad |\omega|=1,
\end{equation}
where $\sigma,e_0>0$ will be small.  The base is
\[
 y_n\ge-B-1,\qquad t\le T_0,\qquad \omega\in\mathbb S^{n-1},
\]
and $a$ decreases through a fixed interval $[a_{\min},a_{\max}]$.
The term $e_0|y'|^2$ makes this base compact. The parameter $a$ decreases
while $(y,\omega)$ remains fixed, as required by the compact-cylinder
propagation lemma.
Take $a_{\min}$ small enough that the target separations exceed
$r_{a_{\min}}(y)$, and $a_{\max}$ large enough that its whole level
satisfies \eqref{cross:coarse-region}.  The bottom base boundary has
$y_n<0$.  Take $T_0$ so large that
$c_0a_{\min}T_0^\sigma>1$; its top boundary also satisfies
\eqref{cross:coarse-region}. Here the superlinear power $1+\sigma$
ensures that $r_a/t=a t^\sigma$ is large on this top face. Thus the
initial $a$-level and both base boundary faces are supplied by the seed.
This known region is upward closed in $a$: since
$\min(x_n,y_n)=y_n+r\min(\omega_n,0)$, its defining
inequality is $y_n<(c_0-\min(\omega_n,0))r$, whose radius coefficient
is positive.

It remains to check the strict extension criterion outside that region.
At the Euclidean metrics use
\[
 \rho=\tfrac12(|x-y|^2-r_a(y)^2),\qquad b=\nabla r_a.
\]
Away from $\omega+b=0$, multiply the vectors in
\eqref{cross:hessian-test} by the common positive factor $r_a$ and
write their real and imaginary parts as
\begin{equation}\label{cross:null-vectors}
 \begin{aligned}
 U_1&=\omega+iv,& |v|=1,\quad v\perp\omega,\\
 U_2&=\frac{\omega+b}{|\omega+b|^2}+iw,
 &w\perp(\omega+b),\quad |w|=|\omega+b|^{-1}.
 \end{aligned}
\end{equation}
The identity $\omega\cdot(U_1-U_2)=b\cdot U_2$ cancels the normal
part of the distance Hessian.  Consequently the test is
\begin{equation}\label{cross:tube-test}
 \big|\pi_{\omega^\perp}(U_1-U_2)\big|^2
 -(r_a\partial^2r_a)(U_2,\overline{U_2}).
\end{equation}
For bounded $b$, outside fixed small neighborhoods of
$b=-\omega$ and $b=-2\omega$, the first term is at least $c|b|^2$.
To see this, a zero forces the real part of $U_2$ to be parallel to
$\omega$, and equality of the projected imaginary lengths gives
$|\omega+b|=1$.  Thus $b=0$ or $b=-2\omega$.
Near $b=0$, the real projected difference controls
$|\pi_{\omega^\perp}b|$ to first order; also
\[
 |\pi_{\omega^\perp}w|
 =1-\omega\cdot b+O(|b|^2),
\]
so the imaginary projected difference controls $|\omega\cdot b|$
to first order.  This proves the quadratic lower bound near zero;
compactness proves it on the remaining range.

Outside \eqref{cross:coarse-region}, $y_n\ge c_0r_a$, so
$r_a/t\le1/c_0$ and
\[
 b=(1+\sigma)(r_a/t)\nabla t
\]
stays bounded.  When $\nabla t$ is close to $e_n$, the two excluded
neighborhoods already lie in \eqref{cross:coarse-region}.  Indeed,
\[
 \frac{x_n}{r_a}
 \le \frac{(1+\sigma)|\nabla t|}{|b|}+\omega_n,
\]
which is arbitrarily small near $b=-\omega$ and negative near
$b=-2\omega$ as $\sigma$ and the gradient error tend to zero.
On the complement $|U_2|$ is bounded and differentiation of
\eqref{cross:tubes} gives
\begin{equation}\label{cross:tube-error}
 r_a|\partial^2r_a|
 \le C(\sigma+T_0e_0)|b|^2.
\end{equation}
Choose the excluded-neighborhood sizes first in terms of $c_0$.
Use the positive lower bound on the remaining bounded $b$-range to
choose $\sigma$ small, then choose the $a$-interval and $T_0$ as above,
and finally $e_0>0$ so small that \eqref{cross:tube-error} is absorbed.
Choosing the transverse compactification last keeps its contribution
to the Hessian test small. These choices also make the gradient error small, since
$|\nabla t-e_n|\le2\sqrt{e_0T_0}$ on the base.
Thus \eqref{cross:tube-test} is strictly positive wherever continuation
is needed.  Both partial gradients are nonzero there.

All these sets are bounded and stay at separation at least
$a_{\min}(M-B-1)^{1+\sigma}>0$.  Smooth convergence of the metrics
therefore preserves the strict tests for sufficiently small $\lambda$.
Lemma~\ref{cross:propagation} continues the kernel through the cylinder,
with uniform estimates and agreement on the seed.  A target pair has
$a=|x-y|/t^{1+\sigma}\ge a_{\min}$; if $a>a_{\max}$ it is already in
the coarse region.  Hence every target pair is covered.

We also identify the extension with both pieces of the full cross
\eqref{cross:mixed-data}. Perform the construction with a larger bound
$B'>B$ and a smaller separation $0<d'<d$, chosen so that for every
$|y|\le B$ the domain
\[
 D_y=B(0,B')\setminus\overline{B(y,d')}
\]
is connected and contains a fixed lower open ball with $x_n<-\eta$.
Choose $B'$ large enough that every excluded closed ball lies strictly
inside $B(0,B')$ and is disjoint from that lower ball.
For fixed $y\in W_\lambda$ in the target range, the continuation
and $G_{1,\lambda}$ solve the same first-variable equation on $D_y$.
On the lower ball the continuation equals $G_{2,\lambda}$, which equals
$G_{1,\lambda}$ because both variables lie in $W_\lambda$.
Unique continuation on $D_y$ proves agreement on the required overlap.
Interchanging the variables proves agreement with the other piece.
Keeping slightly larger target bounds and a slightly smaller separation
before restricting to the prescribed range gives the claimed neighborhood
and its compact-subset estimates.
\end{proof}

The constants in Proposition~\ref{cross:initial} may deteriorate as
$d\downarrow0$.  Its role is to initialize a construction at finitely
many fixed positive scales.  Lemma~\ref{cross:propagation} will also
provide qualitative existence at each positive radius in a shrinking
family; the uniform estimate for that family is established below.
